\specialchapter{书目}

\begin{enumerate}
\def\labelenumi{\arabic{enumi}.}
\item
  F. KLEIN and S. LIE: (a) Sur une certaine famille de courbes et surfaces, C. R. Acad. Sci., 70 (1870), pp.~1222--1226 and 1275--1279 (={[}2{]}, pp.~416--420 and {[}3{]}, vol.~1, pp.~78--85); (b) Über diejenigen ebenen Kurven, welche durch ein geschlossenes System von einfach unendlich vielen vertauschbaren linearen Transformationen in sich übergehen, Math. Ann., 4 (1871), pp.~50--84 (={[}2{]}, pp.~424--459 and {[}3{]}, vol.~1, Abh. XIV, pp.~229--266).
\item
  F. KLEIN, Gesammelte mathematische Abhandlungen, Bd. I, Berlin (Springer), 1921.
\item
  S. LIE, Gesammelte Abhandlungen, 7 vol., Leipzig (Teubner): (a) Über die Reziprozitätsverhältnisse des Reyeschen Komplexes, vol.~I, Abh. V, pp.~68--77 (=Gött. Nach. (1870), pp.~53--66); (b) Über eine Klasse geometrischer Transformationen, vol.~I, Abh. XII, pp.~153--214 (=Christiana For. (1871), pp.~182--245); (c) Über partielle Differentialgleichungen erster Ordnung, vol.~III, Abh. VII, pp.~32--63 (=Christiana For. (1873), pp.~16--51); (d) Theorie der Transformationsgruppen II, vol.~V, Abh. III, pp.~42--75 (=Archiv. Math., I (1876), pp.~152--193); (e) Ueber Gruppen von Transformationen, vol.~V, Abh. I, pp.~1--8 (=Gött. Nachr. (1874), pp.~529--542); (f) Allgemeine Untersuchungen über Differentialgleichungen, die eine kontinuierliche endliche Gruppe gestatten, vol.~VI, Abh. III, pp.~139--223 (=Math. Ann., 25 (1885), pp.~71--151); (g) Beiträge zur allgemeinen Transformationenstheorie, vol.~VI, Abh. V, pp.~230--236 (=Leipziger Ber. (1888), pp.~14--21); (h) Theorie der Transformationsgruppen III, vol.~V, Abh. IV, pp.~78--133 (=Archiv. Math., vol.~III, (1887), pp.~93--165).
\item
  S. LIE and F. ENGEL, Theorie der Transformationsgruppen, 3 vol., Leipzig (Teubner), 1888--1893.
\item
  J. J. SYLVESTER, Collected Mathematical Papers, 4 vols., Cambridge, 1904-1911.
\item
  A. CAYLEY, Collected Mathematical Papers, 13 vols., Cambridge, 1889-1898.
\item
  C. JORDAN, Mémoire sur les groupes de mouvements, Annali di Math., 11 (1868--1869), pp.~167--215 and 332--345 (= Oeuvres, vol.~IV, pp.~231--302).
\item
  F. SCHUR: (a) Zur Theorie der aus Haupteinheiten gebildeten Komplexen, Math. Ann., 33 (1889), pp.~49--60; (b) Neue Begründung der Theorie der endlichen Transformationsgruppen, Math. Ann., 35 (1890), pp.~161--197; (c) Zur Theorie der endlichen Transformationsgruppen, Math. Ann., 38 (1891), pp.~273--286; (d) Über den analytischen Character der eine endliche continuierliche Transformationsgruppe darstellende Funktionen, Math. Ann., 41 (1893), pp.~509--538.
\item
  F. ENGEL: (a) Über die Definitionsgleichung der continuerlichen Transformationsgruppen, Math. Ann., 27 (1886), pp.~1--57; (b) Die Erzeugung der endlichen Transformationen einer projektiven Gruppe durch die infinitesimalen Transformationen der Gruppe, I, Leipzig Ber., 44 (1892), pp.~279--296, II (mit Beiträgen von E. Study), ibid., 45 (1893), pp.~659--696.
\item
  L. MAURER, Über allgemeinere Invarianten-Systeme, Sitzungsber. München, 18 (1888), pp.~103--150.
\item
  W. KILLING, Die Zusammensetzung der stetigen endlichen Transformationsgruppen: (I) Math. Ann., 31 (1888), pp.~252--290; (II) ibid., 33 (1889), pp.~1--48; (III) ibid., 34 (1889), pp.~57--122; (IV) ibid., 36 (1890), pp.~161--189.
\item
  E. CARTAN, Oeuvres complètes, 6 vol., Paris (Gauthier-Villars), 1952--54.
\item
  J. E. CAMPBELL: (a) On a law of combination of operators bearing on the theory of continuous transformation groups, Proc. London Math. Soc., (1) 28 (1897), pp.~381--390; (b) On a law of combination of operators (second paper), ibid., 29 (1898), pp.~14--32.
\item
  H. POINCARÉ, Oeuvres, 11 vol., Paris (Gauthier-Villars), 1916--1956.
\item
  H. F. BAKER, Alternants and continuous groups, Proc. London Math. Soc., (2) 3 (1905), pp.~24--47.
\item
  F. HAUSDORFF, Die symbolische Exponentialformel in der Gruppentheorie, Leipziger Ber., 58 (1906), pp.~19--48.
\item
  A. HURWITZ, Über die Erzeugung der Invarianten durch Integration, Gött, Nachr. (1897), pp.~71--90 (=Math. Werke, vol.~II, pp.~546--564).
\item
  E. E. LEVI, Sulla struttura dei gruppi finiti e continui, Atti Acc. Sci. Torino, 40 (1905), pp.~551--565 (= Opere, vol.~I, pp.~101--115).
\item
  K. HENSEL, Über die arithmetischen Eigenschaften der Zahlen, Jahresber. der D.M.V., 16 (1907), pp.~299--319, 388--393, 474--496.
\item
  I. SCHUR, Neue Anwendungen der Integralrechnung auf Probleme der Invariantentheorie, Sitzungsber. Berlin, 1924, pp.~189--208, 297--321, 346--355.
\item
  H. WEYL, Theorie der Darstellung kontinuierlicher halb-einfacher Gruppen durch lineare Transformationen, I, Math. Zeitschr., 23, (1925),
\end{enumerate}

pp.~271--309; II, ibid., 24 (1926), pp.~328--376; III, ibid., 24 (1926), pp.~377--395 (=Werke, vol.~II, pp.~543--647).

\begin{enumerate}
\def\labelenumi{\arabic{enumi}.}
\setcounter{enumi}{21}
\item
  O. SCHREIER: (a) Abstrakte kontinuierliche Gruppen, Abh. math. Sem. Hamburg, 4 (1926), pp.~15--32; (b) Die Verwandschaft stetiger Gruppen in grossen, ibid., 5 (1927), pp.~233--244.
\item
  J. VON NEUMANN, Zur Theorie der Darstellung kontinuierlicher Gruppen, Sitzungsber. Berlin, 1927, pp.~76--90 (=Collected Works, vol.~I, pp.~134--148).
\item
  P. HALL, A contribution to the theory of groups of prime power order, Proc. London Math. Soc., (3) 4 (1932), pp.~29--95.
\item
  W. MAGNUS: (a) Beziehungen zwischen Gruppen und Idealen in einem speziellen Ring, Math. Ann., 111 (1935), pp.~259--280; (b) Über Beziehungen zwischen höheren Kommutatoren, J. Crelle, 177 (1937), pp.~105--115.
\item
  J. H. C. WHITEHEAD: (a) On the decomposition of an infinitesimal group, Proc. Camb. Phil. Soc., 32 (1936), pp.~229--237 (=Mathematical Works, I, pp.~281--289); (b) Certain equations in the algebra of a semi-simple infinitesimal group, Quart. Journ. of Math., (2) 8 (1937), pp.~220--237 (=Mathematical Works, I, pp.~291--308).
\item
  I. Ado: (a) Note on the representation of finite continuous groups by means of linear substitutions (俄文), Bull. Phys. Math. Soc. Kazan, 7 (1935), pp.~3--43; (b) The representation of Lie algebras by matrices (俄文), Uspehi Mat. Nauk, 2 (1947), pp.~159--173 (英文译本：Amer. Math. Soc. Transl., (1) 9, pp.~308--327).
\item
  N. JACOBSON: (a) Rational methods in the theory of Lie algebras, Ann. of Math., 36 (1935), pp.~875--881; (b) Classes of restricted Lie algebras of characteristic p, II, Duke Math. Journal, 10 (1943), pp.~107--121.
\item
  G. BIRKHOFF: (a) Continuous groups and linear spaces; Rec. Math. Moscou, 1 (1936), pp.~635--642; (b) Representability of Lie algebras and Lie groups by matrices, Ann. of Math., 38 (1937), pp.~526--532.
\item
  E. WITT, Treue Darstellung Lieschen Ringe, J. Creille, 177 (1937), pp.~152--160.
\item
  R. BRAUER, Eine Bedingung für vollständige Reduzibilität von Darstellungen gewöhnlicher und infinitesimaler Gruppen, Math. Zeitschr., 41 (1936), pp.~330--339.
\item
  H. CASIMIR-B. L. VAN DER WAERDEN, Algebraischer Beweis der vollständigen Reduzibilität der Darstellungen halbeinfacher Liescher Gruppen, Math. Ann., 111 (1935), pp.~1--12.
\item
  A. WEIL, Sur les fonctions elliptiques p-adiques, C. R. Acad. Sci., 203 (1935), p.~22.
\item
  E. Lutz, Sur l'équation \(y^{2} = x^{3} - Ax - B\) dans les corps \(p\) -adiques, \(J\) . Crelle, 177 (1937), pp.~237-247.
\item
  C. CHABAUTY, Sur les points rationels des courbes algébriques de genre supérieur à l'unité, C. R. Acad. Sci., 212 (1941), pp.~882--884.
\end{enumerate}

\begin{enumerate}
\def\labelenumi{\arabic{enumi}.}
\setcounter{enumi}{35}
\item
  L. S. PONTRJAGIN, Topological groups, Princeton Univ. Press, 1939.
\item
  R. HOOKE, Linear \(p\) -adic groups and their Lie algebras, Ann. of Math., 43 (1942), pp.~641-655.
\item
  C. CHEVALLEY, Theory of Lie groups, Princeton Univ. Press, 1946.
\item
  E. DYNKIN: (a) Evaluation of the coefficients of the Campbell-Hausdorff formula (俄文), Dokl. Akad. Nauk, 57 (1947), pp.~323-326; (b) Normed Lie algebras and analytic groups (俄文), Uspehi Mat. Nauk, 5 (1950), pp.~135-186 (英文译本：Amer. Math. Soc. Transl., (1) 9, pp.~470-534).
\item
  M. HALL, A basis for free Lie rings and higher commutators in free groups, Proc. Amer. Math. Soc., 1 (1950), pp.~575--581.
\item
  D. MONTGOMERY-L. ZIPPIN, Topological Transformation Groups, New York (Interscience), 1955.
\item
  M. LAZARD: (a) Quelques calculs concernant la formule de Haudorff, Bull. Soc. Math. France, 91 (1963), pp.~435--451; (b) Groupes analytiques p-adiques, Publ. Math. I.H.E.S., no. 26 (1965), pp.~389--603.
\end{enumerate}

\(\mathcal{C}_k\mathfrak{g}\)（g 为李代数）：I.1.6.

ad \(_{g}\) x, ad x（x 为李代数 g 的一个元素）：I.1.2.

\(g_{(K_{1})}\)（g 为李代数）：I.1.9.

\(e^u\) , \(\exp u\)（ \(u\) 为特征为 0 的域上向量空间的幂零自同态）：I.6.8.

\(\mathcal{D}\mathfrak{g}, \mathcal{D}^{\kappa}\mathfrak{g}, \mathcal{C}^{\kappa}\mathfrak{g}\)（g 为李代数）：I.1.5.
% semantic-fragments: legacy-input-seam
\relax{}
